\section{Indefinite spaces of rank at least three}\label{sec:indef}
Let $V$ have indefinite signature and rank $n\ge3$. The HNF
bound and the orientation-union identity continue to hold, while
integral stabilizers may be infinite. We use local isometry types
and embedding spinor norms to count the represented classes. Set
\[
 c_p(e)=\#\{\text{local types of index-}p^e\text{ sublattices of }L_p\},\quad
 c(m)=\prod_{p^e\Vert m}c_p(e),\quad c_p(0)=1 .
\]

\subsection{The spinor norm and represented classes}
The adelic spinor-genus description and Eichler's theorem give
one proper class in each proper spinor genus of an indefinite
rational lattice of rank at least three \cite[\S9.2]{SP}. More explicitly, write $\theta$ for the spinor norm, normalized on a reflection in $v$ by $Q(v)$ modulo squares. On the proper group this convention is independent of a scalar in the quadratic form. For an integral lattice $K$ in this space the proper classes in its genus are parametrized, with $K$ as origin, by the finite elementary abelian group
\begin{equation}\label{eq:spinC}
 C(K)=\mathbb A_f^\times/
 \bigl(\mathbb Q^\times\mathbb A_f^{\times2}
       \prod_p\theta(SO(K_p))\bigr).
\end{equation}
Products of local subgroups are taken in the ideles, with
squareclasses lifted to multiplicative groups. There is no positivity
restriction at the indefinite real place. Strong approximation for the spin group and the spinor-norm
image statements in \cite[\S9.2]{SP} apply to the indefinite space.

The ambient spinor-norm image follows from an indefinite ternary
subspace. Its even Clifford algebra is a quaternion algebra
split at infinity, and conjugation on the ternary space has spinor norm
the reduced norm modulo squares. Every nonzero rational number is a
reduced norm, and the reduced norm is surjective at every finite place.

Improper local isometries introduce no additional genera: choose a
vector $v$ of minimal valuation of nonzero $Q(v)$ in the integral local
lattice. Then $Q(v)$ divides $Q(v+w)-Q(v)-Q(w)$ for every $w$, so
reflection in $v$ is an integral improper isometry, also at two.

For each prime dividing $m$, choose a local type and an embedding
into $L_p$, taking $L_p$ at all other primes. Chinese-remainder
patching gives a reference lattice $K\subset L$ of index $m$.
Different tuples have distinct genera, and all $c(m)$ tuples occur. In \eqref{eq:spinC} let $T(K,L)$ be the set of products of the images under $\theta$ of the sets
\[
 \{g_p\in SO(V_p):g_p K_p\subset L_p\}.
\]
Use integral elements at almost all places. The represented proper classes in this genus are counted by $|T(K,L)|$: patching the local images yields the corresponding spinor class, and every represented class arises from such local embeddings. 
One $g_p$ per local image suffices, since the stabilizer has
trivial image in $C(K)$. For $p\nmid m$, the local image equals $L_p$ and the stabilizer
has trivial image in $C(K)$.

For the ordinary count fix a rational reflection $r$ and choose local improper isometries $j_p\in O(K_p)$, equal to $r$ wherever $r$ is integral except at finitely many places. Put $\delta_K=[\theta((rj_p)_p)]$. This class is independent of the choices of $j_p$; changing $r$ changes the product by a rational proper factor. Reversal acts on $C(K)$ by translation by $\delta_K$, since conjugation by an orthogonal transformation preserves the spinor norm. Thus the number of represented ordinary classes is the cardinality of the image of $T(K,L)$ in $C(K)/\langle\delta_K\rangle$. Summing these formulas over the $c(m)$ tuples gives both actual coefficients at all indices.

\begin{lemma}\label{lem:indlocal}
For $p\nmid 2\det B_L$, $c_p(1)=3$. Each of these types forms a single $SO(L_p)$-orbit of sublattices. For each child $K_p$ its full integral group $SO(K_p)$ has all unit squareclasses as spinor norms. For the isotropic hyperplane type $SO(K_p)$ has all local spinor norms. Moreover $\theta(SO(L_p))$ consists of unit squareclasses.
\end{lemma}
\begin{proof}
A kernel is a hyperplane modulo $p$. Normals are isotropic, of square unit norm, or of nonsquare unit norm. The nonsingular finite quadratic space of rank at least three contains a hyperbolic plane and represents both unit classes. Witt extension gives transitivity on each set of normal lines under the finite orthogonal group. Isometries over $\mathbb F_p$ are products of reflections in nonisotropic vectors; lift each vector to a vector with unit norm to obtain an integral isometry. The determinant can be corrected in a hyperplane stabilizer: in the nonisotropic case use a unit vector in the perpendicular space, and in the isotropic case use the complement of a hyperbolic plane, of rank $n-2\ge1$. Thus each set is a single proper parent orbit.
At a nonisotropic line the kernel has a unimodular rank $n-1$ summand and a unary summand scaled by $p^2$. The latter's unit determinant squareclass distinguishes the two choices, by its intrinsic top Jordan constituent. 

At an isotropic line lift a hyperbolic basis by Hensel lifting and orthogonal splitting. Its kernel replaces the hyperbolic plane $\mathbb H_p$ by $p\mathbb H_p$ (scale of its form) and keeps a unimodular complement. This has a different Jordan decomposition. 

Any unimodular summand of rank at least two represents all unit squareclasses: over $\mathbb F_p$ this follows for a split binary plane by multiplication and for an anisotropic binary plane by the finite-field norm. Lifts preserve the unit squareclass. Products of unit-norm reflections then give every unit spinor class. The same argument in $p\mathbb H_p$ gives all units by products of reflections of norms $pu,pv$; these are integral on the scaled plane. Composing one such reflection with a unit-norm reflection in the complement adds an odd valuation class. 

Finally at a unimodular odd-prime lattice every integral orthogonal transformation is a product of unit-norm reflections. For a unit diagonal basis vector $v$ at least one of $hv\pm v$ has unit norm, since the sum of their norms is $4Q(v)$. Use that reflection, if necessary the reflection in $v$, then induct on its unimodular complement. This proves the last assertion.
\end{proof}
Consequently, for every indefinite parent of rank $n\ge3$,
\begin{equation}\label{eq:indsquarefree}
 a_O(t)=a_{SO}(t)=3^{\omega(t)}\qquad(t\ {\rm squarefree},\ (t,2\det B_L)=1).
\end{equation}
For a fixed tuple, at $p\mid t$ the transporter from the reference child is $SO(L_p)SO(K_p)$ by the single-orbit assertion. Its spinor norm has trivial image in $C(K)$. Elsewhere the same conclusion holds trivially. Thus exactly one proper class is represented per tuple.

\subsection{The leading constant in indefinite signature}
\begin{theorem}\label{th:indmean}
For each of $H=O,SO$, the abscissa is $1$ and
\[
 \sum_{m\le X}a_H(m)\sim \frac{C_L}{2}X(\log X)^2,\qquad
 C_L=\prod_p\left[(1-p^{-1})^3\sum_{e\ge0}\frac{c_p(e)}{p^e}\right]>0,
\]
and in particular $(s-1)^3Z_H(s)\to C_L$ as real $s\downarrow1$.
\end{theorem}
\begin{proof}
Let $S$ contain the primes of $2\det B_L$. Local Jordan splitting bounds $c_p(e)\ll_{n,L}(1+e)^n$ uniformly: the determinant valuation is $v_p(\det B_L)+2e$, and there are at most $(1+v_p(\det B_L)+2e)^n$ ordered scale lists counted with repetition for the $n$ basis positions. At odd $p$ a unimodular constituent is specified by its rank and its unit determinant squareclass. At two there are boundedly many unimodular types of bounded rank, by the dyadic classification \cite{AGM}. These bounds include imprimitive children. 

The multiplicative $c(m)$ therefore has series $\zeta(s)^3 V(s)$ with $V(s)$ an absolutely convergent Dirichlet series for $\Re s>1/2$: at primes outside $S$ the linear Euler coefficient cancels by Lemma~\ref{lem:indlocal}. $V(1)=C_L$ is positive by absolute convergence and positive local values. The elementary estimate $\sum_{j\le x}d_3(j)=x(\log x)^2/2+O(x(1+\log x))$ follows by summing $\lfloor x/(ab)\rfloor$ and the harmonic sum twice. Convolution and dominated convergence give $\sum_{m\le X}c(m)\sim C_L X(\log X)^2/2$.

Write $m=dt$ where $d$ contains all powers at $S$ and all prime powers of exponent at least two outside $S$; $t$ is squarefree and coprime to $S_d=S\cup\{p:p\mid d\}$. At every place outside $S_d$ the local spinor norm of a child stabilizer contains all units. A finite idele may be divided by a positive rational number to make it a unit at all finite places. Its remaining class is generated by the unit squareclasses at $S_d$. Hence the number of proper classes per genus is at most $4\cdot2^{|S_d|}$, with the factor 4 allowing for the dyadic units. This gives a summable majorant after conditioning on $d$, since
\[
 \sum_{d}2^{|S_d|}c(d)/d<\infty .
\]
Indeed all primes outside $S$ dividing such $d$ occur with exponent at least two, and the local polynomial bounds just proved apply. The inequality $\sum_{t\le y}\mu^2(t)3^{\omega(t)}\ll y(1+\log y)^2$ follows already from $3^{\omega(t)}\le d_3(t)$ for squarefree $t$.

Fix $d$ and a tuple of local types at $S_d$. Realize each generator of the product of unit squareclasses at $S_d$ by the residues of odd primes in a reduced arithmetic progression modulo $8\prod_{p\in S_d,\,p>2}p$. CRT supplies the progressions, after excluding finitely many primes. If for each of these finitely many generators some such $q\mid t$ carries the isotropic local type, the proper genus has only one class. In fact the idele supported at $q$ with uniformizer $q$ is trivial in \eqref{eq:spinC} by Lemma~\ref{lem:indlocal}. Dividing it by rational $q$ kills the prescribed unit tuple at $S_d$ and leaves only units at other primes, where those units are in the stabilizer image. Thus all generators vanish. Both actual counts give one class for this tuple.

For a particular generator choose a finite subset $F$ of its prime progression. Among squarefree $t$ coprime to $S_d$, giving three labels to every prime, the asymptotic proportion of labelled tuples using no isotropic label at any $q\in F$ equals
\[
 \prod_{q\in F}\frac{1+2/q}{1+3/q}.
\]
At these primes the local factor $1+3q^{-s}$ is replaced by
$1+2q^{-s}$. Extracting $\zeta(s)^3$ from the two products and
applying the preceding convolution estimate gives the proportion. The proportions tend to zero as $F$ exhausts the progression, by the divergence of the reciprocals of primes in reduced residue classes. A finite union over generators shows that all exceptional labelled tuples at fixed $d$ contribute $o_d(X(\log X)^2)$, even after multiplying by the fixed genus-size bound. Finally truncate the summable majorant in $d$ to pass to the full sum. Thus the excess of either actual count over the genus count is $o(X(\log X)^2)$. Partial summation gives the stated real boundary limit, divergence at $1$, and absolute convergence to its right.
\end{proof}
